%%%%%%%%%%%%%%%%%%%%%%%%
%
% $Autor: Wings $
% $Project:  $
% $Filename: OledSSD1306
% $Short description: In diesem Abschnitt wird das OLED-Display genauer beschrieben. Wie dieses aufgebaut ist und was benötigt wird, um dieses in ein Projekt einzubinden.
%
%%%%%%%%%%%%

\chapter{OLED-Display}\label{sec:OLEDDisplay}
\section{Allgemein}
Ein \ac{oled} stellt eine besondere Form der LED-Technologie dar, bei der jedes Pixel aus eigenständig leuchtenden organischen Leuchtdioden besteht. Im Gegensatz zu LCDs mit LED-Hintergrundbeleuchtung benötigen OLEDs keine Hintergrundbeleuchtung oder LC-Zellen, da die einzelnen Pixel selbst Licht erzeugen. Dies ermöglicht nicht nur eine besonders flache Bauweise, sondern auch eine exakte Steuerung einzelner Bildpunkte – inklusive vollständigem Abschalten für tiefstes Schwarz und hohen Kontrast.

Ein wesentlicher Vorteil gegenüber LCD-Technologien ist die deutlich schnellere Reaktionszeit: Helligkeitsänderungen erfolgen bei OLEDs in unter einer Mikrosekunde. Auch der Farbraum und Kontrastumfang sind bei OLED-Bildschirmen erweitert.
In der Praxis kommen verschiedene OLED-Technologien zum Einsatz, insbesondere RGB-SBS-AMOLED (mit drei nebeneinander angeordneten Farb-OLEDs pro Pixel) und WOLED (weiße OLEDs mit nachgeschalteten Farbfiltern). OLED-Elemente werden meist im Dünnschichtverfahren hergestellt, was die Produktion vereinfacht.\cite{Stotz:2019}



S\section{OLED-Display mit SSD1306 Controller}
Das \ac{oled}-Display mit dem Controller SSD1306 dient der Ausgabe von Messwerten des Arduinos. Es besitzt Abmessungen von 27 x 27 mm und überzeugt durch seinen hohen Kontrast sowie gute Lesbarkeit. Die Darstellung erfolgt in blauer Farbe. Angesteuert wird das Display über die I\textsuperscript{2}C-Schnittstelle (Pins VCC, GND, SCL, SDA) und benötigt eine Betriebsspannung von 3,3 V bis 5 V. Der zulässige Betriebstemperaturbereich liegt zwischen -30 °C und +70 °C. Für die Ansteuerung kommt die Adafruit\_GFX-Bibliothek zum Einsatz, welche eine Vielzahl von Grafikfunktionen für \ac{oled}-Displays mit SSD1306-Controller bereitstellt.\cite{Velleman:WPI438}

\begin{longtable}[h]{|p{5cm}|p{8cm}|}
	\caption{Spezifikationen OLED-Display}
	\label{tab:specs_oled}\\
	\hline
	%Artikel Identifikation
	\multicolumn{2}{|l|}{\rule{0pt}{15pt}\textbf{Identifikation}} \\ \hline
	Artikel-Nr.& ARD OLED 0.96 \\ \hline
	EAN/GTIN & 5410329729738 \\ \hline
	Hst.-Teile-Nr. & WPI438 \\ \hline
	% Hauptmerkmale
	\multicolumn{2}{|l|}{\rule{0pt}{15pt}\textbf{Hauptmerkmale}} \\ \hline
	Modell & OLED-Display \\ \hline
	Abmessungen & 26,7 x 31,26 \\ \hline
	Gewicht & 1,54 g \\ \hline
	Lieferumfang & WPI438 \\ \hline
	
	% Weitere Spezifikationen
	\multicolumn{2}{|l|}{\rule{0pt}{15pt}\textbf{Weitere Spezifikationen}} \\ \hline
	Spannungsversorgung & 3.3 V bis 5 V \\ \hline
	Auflösung & 128 x 64 Pixel \\ \hline
	Hintergrundfarbe & Schwarz \\ \hline
	Betrachtungswinkel & > 160° \\ \hline
	Anzeigefarbe & White \\ \hline
	Interface & I\textsuperscript{2}C \\ \hline
	Controller-Chip & SSD1306 \\ \hline
	
\end{longtable}


\subsection{Aufbau und Funktionsweise}
In der Abbildung \ref{fig:OLEDsketch} wird der Aufbau des OLED-Displays gezeigt. Über die vier Pins wird die Stromversorgung und Logikspannung des OLED-Displays mit dem Mikrocontroller verbunden.

\begin{center}
	\begin{circuitikz}[scale=0.6, transform shape]
	\centering
	\tikzstyle{every node}=[font=\small]
	
	\draw [ fill={rgb,255:red,24; green,9; blue,78} , line width=1pt ] (2.5,18.5) rectangle (11.25,7.25);
	
	\node [color=white, rotate=90] at (3.5,13.25) {SCK};
	
	\draw [ fill=black , line width=1pt ] (4,18) rectangle (9.75,7.75);
	\draw [ fill={rgb,255:red,40; green,10; blue,123} , line width=1pt ] (4.25,17.75) rectangle (9,8);
	
	\draw [ fill={rgb,255:red,188; green,184; blue,184} ] (10.5,17.75) circle (0.5cm);
	\draw [ dashed] (10.5,17.75) circle (0.25cm);
	
	\draw [ fill={rgb,255:red,188; green,184; blue,184} ] (3.25,8) circle (0.5cm);
	\draw [ fill={rgb,255:red,188; green,184; blue,184} ] (3.25,17.75) circle (0.5cm);
	\draw [ fill={rgb,255:red,188; green,184; blue,184} ] (10.5,8) circle (0.5cm);
	
	\draw [ dashed] (10.5,8) circle (0.25cm);
	\draw [ dashed] (3.25,17.75) circle (0.25cm);
	\draw [ dashed] (3.25,8) circle (0.25cm);
	
	\draw [ fill={rgb,255:red,186; green,186; blue,186} , rounded corners = 9.6] (10,14.75) rectangle (11.25,11);
	\draw [ fill={rgb,255:red,186; green,186; blue,186}] (10.75,14.75) rectangle (11.25,11);
	
	\draw [ color=red , fill={rgb,255:red,255; green,184; blue,184}, dashed] (2.75,14.25) rectangle  (3.25,13.75);
	\draw [ color=red , fill={rgb,255:red,255; green,184; blue,184}, dashed] (2.75,13.5) rectangle  (3.25,13);
	\draw [ color=red , fill={rgb,255:red,255; green,184; blue,184}, dashed] (2.75,12.75) rectangle  (3.25,12.25);
	\draw [ color=red , fill={rgb,255:red,255; green,184; blue,184}, dashed] (2.75,12) rectangle  (3.25,11.5);
	
	\draw [ color=red , fill={rgb,255:red,255; green,184; blue,184}] (3,14) circle (0.25cm);
	\draw [ color=red , fill={rgb,255:red,255; green,184; blue,184}] (3,13.25) circle (0.25cm);
	\draw [ color=red , fill={rgb,255:red,255; green,184; blue,184}] (3,12.5) circle (0.25cm);
	\draw [ color=red , fill={rgb,255:red,255; green,184; blue,184}] (3,11.75) circle (0.25cm);
	
	\draw [ fill={rgb,255:red,208; green,184; blue,27}] (10,14.25) rectangle (10.5,11.5);
	
	% Pins
	\node [color=white, rotate=90] at (3.5,14) {SDA};
	\node [color=white, rotate=90] at (3.5,13.25) {SCL};
	\node [color=white, rotate=90] at (3.5,12.5) {VDD};
	\node [color=white, rotate=90] at (3.5,11.75) {GND};
	
	\node [color=white, rotate=90] at (3,14.5) {4};
	\node [color=white, rotate=90] at (3,11.25) {1};
	
	\draw [ fill={rgb,255:red,130; green,130; blue,130} ] (10,14) rectangle (10.25,11.75);
	
	% Beschriftung außen
	\node at (-1,14) {I²C Serial Data Line - SDA};
	\node at (-1,13.25) {I²C Serial Clock Line - SCL};
	\node at (-1,12.5) {Power Supply - VDD};
	\node at (-1,11.75) {Ground - GND};
	
	\draw[-, thick] (12,10) -- (8,10);
	\node[right] at (12,10) {OLED-Display 128 x 64 Pixel};
	
\end{circuitikz}

	\captionof{figure}[OLED-Display]{OLED-Display nach}\label{fig:OLEDsketch}
	\href{../../Datenblaetter/Oled_VMA438_A4V01.pdf}{Oled\_VMA438\_A4V01.pdf}
	
\end{center}

Bei der Ansteuerung eines OLED-Displays über das I\textsuperscript{2}C-Protokoll übernehmen die beiden Leitungen \ac{sda} und \ac{scl} den seriellen Datenaustausch.

\begin{itemize}
	\item SDA überträgt die eigentlichen Datenbits (z.B. Befehle oder Bildinhalte).
	\item SCL gibt den Takt vor, der bestimmt, wann ein Bit gelesen oder geschrieben wird.
\end{itemize}

Das I\textsuperscript{2}C-Protokoll ermöglicht eine einfache, bidirektionale Kommunikation über nur zwei Leitungen, wobei mehrere Geräte über individuelle Adressen angesprochen werden können.\cite{Velleman:WPI438UserManual}


\section{Installation}
Zur Nutzung des OLED-Displays muss dieses mit dem Arduino nach Tabelle \ref{tab:oled_pinbelegung} angeschlossen werden.  
\begin{longtable}[h]{|p{5cm}|p{8cm}|}
	\caption{Pinbelegung des OLED-Displays}
	\label{tab:oled_pinbelegung} \\
	\hline
	\textbf{OLED-Display Pin} & \textbf{Arduino Pin} \\ \hline
	SDA & A4 \\ \hline
	SCL & A5 \\ \hline
	VCC & 5V out \\ \hline
	GND & GND \\ \hline
\end{longtable}




\section{Bibliotheken}
Für die Verwendung des OLED-Displays sind die folgenden aufgelisteten Bibliotheken notwendig \ref{tab:oled_libs}. Um eine Bibliothek in ein Programm zu integrieren muss diese in der Arduino IDE heruntergeladen und installiert sein und anschließend im Programmcode mit dem folgenden Befehl eingebunden werden: \PYTHON{\#include<name.h>} \newline


\section{Verwendete Bibliotheken}
\begin{center}
	\captionof{table}{Verwendete Bibliotheken zum Testen des OLED-Displays}
	\label{tab:oled_libs}
	\begin{tabular}{|p{5cm}|p{9cm}|}
		\hline
		\textbf{Bibliothek} & \textbf{Beschreibung} \\
		\hline
		\textbf{Wire.h} &  \textbf{Version:} Arduino-Standardbibliothek \newline
		Funktion: Dient der Kommunikation über den I\textsuperscript{2}C-Bus zwischen Mikrocontroller und OLED-Display. \\
		\hline
		\textbf{Adafruit\_GFX.h} & \textbf{Version:} 1.12.6 \newline
		Grafikbibliothek zur Darstellung von Text, Linien, Rechtecken und weiteren primitiven Elementen auf grafischen Displays. \\
		\hline
		\textbf{Adafruit\_SSD1306.h} & \textbf{Version:} 2.5.16 \newline
		Treiber für SSD1306-basierte OLED-Displays mit Unterstützung für I\textsuperscript{2}C. \\
		\hline
	\end{tabular}
\end{center}

\subsection{Wire.h}
Die Bibliothek Wire.h ist eine Standardbibliothek für Arduino-Plattformen, welche Funktionen und Methoden für die Kommunikation zur Verfügung stellt. Mit Hilfe der Schnittstelle I\textsuperscript{2}C ermöglicht die Bibliothek die Kommunikation zwischen einem Arduino und einem anderen I\textsuperscript{2}C-fähigen Gerät wie z.B. Sensoren oder Displays. I\textsuperscript{2}C ist ein Kommunikationsbus, der im Vergleich zu seriellen Schnittstellen den Vorteil hat, dass er mit mehr als zwei Geräten kommunizieren kann. Ein I\textsuperscript{2}C-Bus benötigt zwei Leitungen: SCL für ein Taktsignal und SDA für Daten.\cite{arduinoLibraries:2025}

\subsection{Adafruit\_GFX.h}
Die Adafruit\_GFX-Bibliothek bietet grundlegende Grafikfunktionen für viele verschiedene LCD- und OLED-Displays sowie LED-Matrizen von Adafruit. Sie sorgt dafür, dass Programme leicht zwischen unterschiedlichen Displays angepasst werden können. Neue Funktionen und Verbesserungen gelten automatisch für alle unterstützten Displays.
Die Bibliothek wird immer zusammen mit einer weiteren Bibliothek verwendet, die speziell für das jeweilige Display gedacht ist, für unser spezifisches OLED-Display, welches den Controller-Chip SSD1306 verwendet, muss die Adafruit\_SSD1306.h Bibliothek ebenfalls mit eingebunden werden.\cite{AdafruitGFX:2023}

\subsubsection{Adafruit\_SSD1306.h}
Die Adafruit\_SSD1306 Bibliothek ist speziell für OLED-Displays mit dem SSD1306-Controller entwickelt worden. Sie arbeitet zusammen mit der Adafruit\_GFX-Bibliothek und ermöglicht die Ansteuerung und Darstellung von Grafiken und Text auf dem Display.\cite{AdafruitGFX:2023}


%\subsection{Kalibrierung}


\section{Codebeispiele}

Zur Überprüfung des OLED-Displays wurden zwei Testprogramme erstellt. 
Der erste Test ist ein einfacher Funktionstest, bei dem alle Pixel des Displays im Abstand von zwei Sekunden ein- und ausgeschaltet werden. 
Der zweite Test zeigt eine Fortschrittsanzeige, die durch einen Taster ausgelöst wird. 
Dadurch können sowohl die grundlegende Funktion des Displays als auch die grafische Darstellung und die Verarbeitung eines Eingangssignals überprüft werden.

\subsection{Verwendete Funktionen}

In den Testprogrammen für das OLED-Display werden verschiedene Funktionen zur Initialisierung, Ausgabe und Steuerung des Displays verwendet:

\begin{itemize}
	\item \PYTHON{Serial.begin()}: Startet die serielle Kommunikation mit 9600 Baud.
	
	\item \PYTHON{Serial.println()}: Gibt eine Textzeile im Serial Monitor aus.
	
	\item \PYTHON{pinMode()}: Konfiguriert einen Pin als Ein- oder Ausgang.
	
	\item \PYTHON{digitalRead()}: Liest den aktuellen Zustand eines digitalen Eingangs.
	
	\item \PYTHON{display.begin()}: Initialisiert das OLED-Display mit angegebener I\textsuperscript{2}C-Adresse.
	
	\item \PYTHON{display.clearDisplay()}: Löscht alle Inhalte auf dem Display.
	
	\item \PYTHON{display.fillScreen()}: Füllt den gesamten Anzeigebereich mit einer angegebenen Farbe.
	
	\item \PYTHON{display.setTextSize()}: Setzt die Schriftgröße für die Textausgabe.
	
	\item \PYTHON{display.setTextColor()}: Legt die Farbe der Textausgabe fest.
	
	\item \PYTHON{display.setCursor()}: Positioniert den Textcursor an einer bestimmten Stelle.
	
	\item \PYTHON{display.print()}: Gibt Text oder Werte auf dem Display aus.
	
	\item \PYTHON{display.drawRect()}: Zeichnet den Rahmen eines Rechtecks auf dem Display.
	
	\item \PYTHON{display.fillRect()}: Zeichnet ein ausgefülltes Rechteck auf dem Display.
	
	\item \PYTHON{display.display()}: Überträgt die Zeichenpufferdaten auf das Display.
	
	\item \PYTHON{delay()}: Erzeugt eine zeitliche Verzögerung im Programmablauf.
\end{itemize}

\subsection{Globale Variablen und Parameter}

In den Testprogrammen werden folgende globale Variablen und Parameter verwendet:

\begin{itemize}
	\item \PYTHON{SCREEN\_WIDTH}: Definiert die Breite des Displays in Pixeln mit 128 Pixeln.
	
	\item \PYTHON{SCREEN\_HEIGHT}: Definiert die Höhe des Displays in Pixeln mit 64 Pixeln.
	
	\item \PYTHON{OLED\_RESET}: Legt den Reset-Pin des OLED-Displays fest. Der Wert -1 bedeutet, dass kein separater Reset-Pin verwendet wird.
	
	\item \PYTHON{OLED\_ADDRESS}: Definiert die I\textsuperscript{2}C-Adresse des OLED-Displays.
	
	\item \PYTHON{BUTTON\_PIN}: Definiert den digitalen Eingangspin für den Taster.
	
	\item \PYTHON{display}: Erstellt ein Objekt zur Steuerung des OLED-Displays mit dem Adafruit-SSD1306-Treiber.
	
	\item \PYTHON{SSD1306\_SWITCHCAPVCC}: Legt die Spannungsversorgung des Displays über die interne Ladungspumpe des SSD1306-Controllers fest.
	
	\item \PYTHON{SSD1306\_WHITE}: Gibt die Farbe für eingeschaltete Pixel des OLED-Displays an.
	
	\item \PYTHON{INPUT\_PULLUP}: Aktiviert den internen Pull-up-Widerstand des Mikrocontrollers für den Tastereingang.
	
	\item \PYTHON{LOW}: Beschreibt den Zustand eines digitalen Eingangs mit niedrigem Pegel. Beim verwendeten Pull-up-Eingang bedeutet dieser Zustand, dass der Taster gedrückt ist.
	
	\item \PYTHON{i}: Laufvariable zur schrittweisen Erhöhung des Fortschritts.
	
	\item \PYTHON{percent}: Speichert den aktuellen Fortschritt in Prozent.
	
	\item \PYTHON{barWidth}: Speichert die berechnete Breite des gefüllten Fortschrittsbalkens.
\end{itemize}


\subsection{Funktionstest des OLED-Displays}

\subsubsection{Ziel des Tests}

Das Programm \textit{Funktionstest} dient zur Überprüfung des OLED-Displays auf fehlerhafte oder tote Pixel. Dazu werden alle Pixel des Displays im Abstand von zwei Sekunden vollständig ein- und ausgeschaltet. Durch dieses wiederholte Blinken kann kontrolliert werden, ob alle Bildpunkte korrekt funktionieren.

\subsubsection{Anschluss des OLED-Displays}

Das OLED-Display wird über die I\textsuperscript{2}C-Schnittstelle angeschlossen. Die Spannungsversorgung erfolgt über VCC und GND, während SDA für die Datenübertragung und SCL für das Taktsignal verwendet wird.

\begin{itemize}
	\item VCC an 3,3 V oder 5 V
	\item GND an GND
	\item SDA an A4
	\item SCL an A5
\end{itemize}

\subsubsection{Programmablauf}

Im Programm wird das Display zunächst initialisiert und anschließend geleert. Danach werden in der Funktion \PYTHON{loop()} alle Pixel mit \PYTHON{display.fillScreen()} eingeschaltet. Nach einer Wartezeit von zwei Sekunden wird das Display mit \PYTHON{display.clearDisplay()} wieder ausgeschaltet. Dieser Ablauf wiederholt sich dauerhaft, wodurch mögliche Pixelfehler sichtbar werden.

{
	\captionof{code}{Funktionstest für das OLED-Display}
	\label{Code:Test:File:TestOLEDDisplay}
	\ArduinoExternal{}{../../Code/OledFunktionsTest/OledFunktionsTest.ino}  
}


\subsection{Fortschrittsanzeige auf dem OLED-Display}

\subsubsection{Ziel des Tests}

In diesem Test wird überprüft, ob durch Betätigung eines Tasters eine visuelle Rückmeldung auf dem OLED-Display ausgegeben werden kann. Beim Drücken des Tasters wird ein Fortschrittsbalken von 0 \% bis 100 \% in zehn Schritten angezeigt. Dadurch kann sowohl die Funktion des Tasters als auch die grafische Ausgabe auf dem OLED-Display überprüft werden.

\subsubsection{Anschluss von Taster und OLED-Display}

Der Taster wird an den digitalen Pin D2 des Mikrocontrollers angeschlossen. Ein Anschluss des Tasters wird mit D2 verbunden, der andere mit GND. Dabei wird der interne Pull-up-Widerstand des Mikrocontrollers verwendet.

Das OLED-Display wird ebenfalls über die I\textsuperscript{2}C-Schnittstelle angeschlossen. Die Pins werden wie folgt verbunden:

\begin{itemize}
	\item VCC an 3,3 V oder 5 V
	\item GND an GND
	\item SDA an A4
	\item SCL an A5
	\item Taster an D2 und GND
\end{itemize}

\subsubsection{Programmablauf}

In der Funktion \PYTHON{loop()} wird der Zustand des Tasters abgefragt. Da \PYTHON{INPUT\_PULLUP} verwendet wird, entspricht der Zustand \PYTHON{LOW} einem gedrückten Taster. Wird der Taster betätigt, wird der Fortschrittsbalken schrittweise gefüllt. Zusätzlich wird der aktuelle Fortschritt als Prozentwert auf dem OLED-Display angezeigt. Zwischen den einzelnen Schritten wird eine kurze Verzögerung eingefügt, damit der Fortschritt sichtbar dargestellt wird.

{
	\captionof{code}{Fortschrittsbalken auf dem OLED-Display durch Tastendruck}
	\label{Code:Test:File:OLEDProgressBar}
	\ArduinoExternal{}{../../Code/OLEDProgressBar/OLEDProgressBar.ino}
}


